% !TEX root = ../main.tex
% Modelo de datos de la coleccion.
%
% Campos y recuentos COMPROBADOS el 29/08/2026 contra los ficheros reales
% data/grados.json (837 items: 528 asignatura, 288 guia, 12 grado, 8 salidas y
% 1 procedencia) y data/chunks.json (1.500 entradas: 1.499 fragmentos y 1
% procedencia). No se escriben a mano: salen de leer los ficheros.
%
% Respecto a la version anterior de esta figura se corrigen dos cosas:
%
%   1. Dibujaba CUATRO origenes de fragmento y hay SIETE. Faltaban mencion
%      (24), ficha_titulacion (12) y catalogo (3), es decir 39 fragmentos que
%      existen en la coleccion y que la figura negaba.
%   2. La caja de procedencia juntaba en una sola los campos de DOS cabeceras
%      distintas. La de grados.json trae dos campos y la de chunks.json trae
%      cinco, y no son los mismos.
%
% El corte en dos bandas es el de los dos ficheros, que es tambien el que
% separa lo extraido de lo indexado. De los 1.499 fragmentos, 1.318 llevan
% texto publicado por la Escuela (guia y salidas), 95 los compone el
% fragmentador contando y agrupando (plan, ficha, mencion y catalogo) y 86 son
% el aviso de que no hay guia publicada.
%
% `ects` es una cadena y no un decimal porque la fuente publica "6" y una
% asignatura lo trae vacio; imputarlo seria inventar un dato que no esta.
%
% Las cajas se encadenan con `below=of` y todo lo demas se coloca RELATIVO a
% ellas. La version anterior daba las alturas a mano y se solapaban: cada caja
% mide lo que miden sus campos, y eso no se puede estimar de antemano.
\begin{figure}[H]
  \centering
  \resizebox{\textwidth}{!}{%
    \begin{tikzpicture}[
      clase/.style={rectangle split, rectangle split parts=2,
        draw=diagNodoBorde, fill=diagNodo, align=left,
        font=\footnotesize, inner sep=5pt, text width=4.6cm},
      claseFoco/.style={clase, draw=diagFocoBorde, fill=diagFocoFondo,
        very thick, text width=4.9cm},
      claseAux/.style={clase, draw=diagNodoBorde, fill=white,
        densely dashed, text width=9.6cm},
      mult/.style={font=\scriptsize, inner sep=1.5pt},
      % La máscara blanca cae siempre fuera de la banda amarilla; por eso se
      % separan las etiquetas 0,65 cm del borde de su caja.
      rotulo/.style={font=\scriptsize, fill=white, inner sep=1.5pt,
        align=left, anchor=south west},
      ]

      % ── Lo que el rastreador extrae de la web ──────────────────────────────
      \node[clase] (grado) {\textbf{Grado}
        \nodepart{two}
        nombre \textcolor{diagNodoBorde}{\texttt{str}}\\
        es\_doble\_grado \textcolor{diagNodoBorde}{\texttt{bool}}\\
        url\_asignaturas \textcolor{diagNodoBorde}{\texttt{str}}\\
        url\_salidas \textcolor{diagNodoBorde}{\texttt{str}}};

      \node[clase, below=1.1cm of grado] (asig) {\textbf{Asignatura}
        \nodepart{two}
        grado \textcolor{diagNodoBorde}{\texttt{str}}\\
        codigo \textcolor{diagNodoBorde}{\texttt{str}}\\
        nombre \textcolor{diagNodoBorde}{\texttt{str}}\\
        tipo\_asignatura \textcolor{diagNodoBorde}{\texttt{str}}\\
        ects \textcolor{diagNodoBorde}{\texttt{str}}\\
        curso \textcolor{diagNodoBorde}{\texttt{str}}\\
        cuatrimestre \textcolor{diagNodoBorde}{\texttt{str}}\\
        ofertada \textcolor{diagNodoBorde}{\texttt{bool}}\\
        tiene\_guia \textcolor{diagNodoBorde}{\texttt{bool}}\\
        menciones \textcolor{diagNodoBorde}{\texttt{list}}\\
        url\_guia \textcolor{diagNodoBorde}{\texttt{str}}};

      \node[clase, below=1.1cm of asig] (guia) {\textbf{Guía}
        \nodepart{two}
        grado \textcolor{diagNodoBorde}{\texttt{str}}\\
        codigo \textcolor{diagNodoBorde}{\texttt{str}}\\
        nombre \textcolor{diagNodoBorde}{\texttt{str}}\\
        resumen \textcolor{diagNodoBorde}{\texttt{str}}\\
        temario \textcolor{diagNodoBorde}{\texttt{str}}\\
        fallback \textcolor{diagNodoBorde}{\texttt{bool}}\\
        curso \textcolor{diagNodoBorde}{\texttt{str}}\\
        formato \textcolor{diagNodoBorde}{\texttt{str}}};

      \node[clase, below=1.1cm of guia] (salidas) {\textbf{Salidas}
        \nodepart{two}
        grado \textcolor{diagNodoBorde}{\texttt{str}}\\
        texto \textcolor{diagNodoBorde}{\texttt{str}}};

      % ── Lo que se indexa ───────────────────────────────────────────────────
      % El fragmento se centra a la altura del hueco entre Asignatura y Guía,
      % que es el centro vertical de la columna. El punto medio se saca con un
      % `midway` sobre un camino invisible: no hace falta la librería `calc`.
      \path (asig.south) -- (guia.north) node[midway] (medio) {};
      \node[claseFoco] (frag) at (12,0 |- medio) {\textbf{Fragmento}
        \nodepart{two}
        origen \textcolor{diagNodoBorde}{\texttt{str}}\\
        grados \textcolor{diagNodoBorde}{\texttt{list}}\\
        codigos \textcolor{diagNodoBorde}{\texttt{list}}\\
        nombre \textcolor{diagNodoBorde}{\texttt{str}}\\
        tipo\_asignatura \textcolor{diagNodoBorde}{\texttt{str}}\\
        curso \textcolor{diagNodoBorde}{\texttt{str}}\\
        texto \textcolor{diagNodoBorde}{\texttt{str}}\\
        chunk\_index \textcolor{diagNodoBorde}{\texttt{int}}\\
        total\_chunks \textcolor{diagNodoBorde}{\texttt{int}}};

      % Los cuatro puntos de llegada, separados para que ninguna flecha tape
      % a otra al entrar.
      \coordinate (e1) at ([yshift=1.5cm]frag.west);
      \coordinate (e2) at ([yshift=0.5cm]frag.west);
      \coordinate (e3) at ([yshift=-0.5cm]frag.west);
      \coordinate (e4) at ([yshift=-1.5cm]frag.west);

      % ── Las dos bandas son los dos ficheros ────────────────────────────────
      % La coordenada suelta mete dentro de la banda el rodeo de Grado a
      % Salidas, que si no quedaría fuera del fichero al que pertenece.
      \begin{scope}[on background layer]
        \node[grupoDiag, fit={(grado)(asig)(guia)(salidas)(-3.6,0)},
        label={[tituloGrupo, anchor=south]north:Extraído de la web: \texttt{grados.json}}] {};
        \node[grupoDiag, fit=(frag),
        label={[tituloGrupo, anchor=south]north:Indexado: \texttt{chunks.json}}] {};
      \end{scope}

      % ── Relaciones entre lo extraído ───────────────────────────────────────
      \draw[flechaDiag] (grado.south) --
      node[mult, right]{1 $\rightarrow$ *} (asig.north);
      \draw[flechaDiag] (asig.south) --
      node[mult, right]{1 $\rightarrow$ 0..1} (guia.north);
      % Rodea la columna por la izquierda: por dentro pasaría por detrás de
      % Asignatura y de Guía, que no son ni origen ni destino.
      \draw[flechaDiag, rounded corners=6pt]
      (grado.west) -- (-3.2,0 |- grado.west) -- (-3.2,0 |- salidas.west)
      -- (salidas.west);
      % Con desplazamiento positivo caían DENTRO de la caja, encima de un campo.
      \node[mult, anchor=south east] at ([xshift=-0.12cm, yshift=0.06cm]grado.west) {1};
      \node[mult, anchor=north east] at ([xshift=-0.12cm, yshift=-0.06cm]salidas.west) {0..1};

      % ── Las siete procedencias de un fragmento ─────────────────────────────
      % Las verticales van escalonadas a propósito: con cualquier otro orden
      % la horizontal de una flecha cruzaría la vertical de otra.
      \draw[flechaDiag, rounded corners=6pt]
      (grado.east) -- (7.2,0 |- grado.east) -- (7.2,0 |- e1) -- (e1);
      \node[rotulo] at ([xshift=0.65cm, yshift=0.1cm]grado.east)
      {plan\_de\_estudios 56\\ficha\_titulacion 12\\mencion 24 · catalogo 3};

      \draw[flechaDiag, rounded corners=6pt]
      (asig.east) -- (6.4,0 |- asig.east) -- (6.4,0 |- e2) -- (e2);
      \node[rotulo] at ([xshift=0.65cm, yshift=0.1cm]asig.east)
      {asignatura\_sin\_guia 86};

      \draw[flechaDiag, rounded corners=6pt]
      (guia.east) -- (7.8,0 |- guia.east) -- (7.8,0 |- e3) -- (e3);
      \node[rotulo] at ([xshift=0.65cm, yshift=0.1cm]guia.east) {guia 1.296};

      \draw[flechaDiag, rounded corners=6pt]
      (salidas.east) -- (8.5,0 |- salidas.east) -- (8.5,0 |- e4) -- (e4);
      \node[rotulo] at ([xshift=0.65cm, yshift=0.1cm]salidas.east) {salidas 22};

      % ── Dos apuntes que el esquema por sí solo no dice ─────────────────────
      \node[font=\scriptsize, align=left, anchor=north west]
      at ([xshift=-0.35cm, yshift=-1.1cm]frag.south west)
      {\texttt{grados} y \texttt{codigos} son listas paralelas: una guía\\
        compartida se imparte en varias titulaciones.};

      \node[claseAux, anchor=north] at (8,0 |- salidas.south)
      {\textbf{Procedencia} --- encabeza cada fichero y no se indexa
        \nodepart{two}
        \texttt{grados.json} \textcolor{diagNodoBorde}{\texttt{fecha\_extraccion, origen}}\\
        \texttt{chunks.json} \textcolor{diagNodoBorde}{\texttt{fecha\_extraccion,
            fecha\_troceado, cursos, tamanos, guias\_sin\_curso}}\\
        Son dos cabeceras distintas, no una: cada fichero anota lo suyo.};
    \end{tikzpicture}%
  }
  \caption[Modelo de datos de la colección]{\textcolor{borrador2}{Modelo de datos
      de la colección. A la izquierda, lo que el rastreador extrae de la web; a
      la derecha, el fragmento que se indexa. Cada flecha lleva escrito qué valor
      toma el campo \texttt{origen} del fragmento y cuántos produce. Fuente:
      Elaboración propia.}}
  \label{fig:modelo_datos}
\end{figure}
